% PRIMARY-SOURCE-EXEMPT: reason=primary research paper; citations in bibliography: Wall (1960) doi:10.1080/00029890.1960.11989541, Fibonacci (1202) Liber Abaci, Sierpinski (1962) ISBN 978-0-486-43293-4
\documentclass[12pt]{article}
\usepackage{amsmath, amssymb, amsthm}
\usepackage{geometry}
\usepackage{graphicx}
\usepackage{booktabs}
\usepackage{array}
\usepackage{multirow}
\usepackage{xcolor}
\usepackage{hyperref}

\geometry{margin=1in}

\newtheorem{theorem}{Theorem}
\newtheorem{lemma}[theorem]{Lemma}
\newtheorem{proposition}[theorem]{Proposition}
\newtheorem{corollary}[theorem]{Corollary}
\newtheorem{conjecture}[theorem]{Conjecture}
\newtheorem{definition}{Definition}
\newtheorem{remark}{Remark}
\newtheorem{example}{Example}

\title{A Complete Classification of Pythagorean Triples\\via Generalized Fibonacci Sequences and Pisano Periods}

\author{[Your Name]\\[Your Affiliation]}

\date{\today}

\begin{document}

\maketitle

\begin{abstract}
We present a comprehensive hierarchical classification of Pythagorean triples based on three nested taxonomic layers. The foundation is a partition into five disjoint families---Fibonacci, Lucas, Phibonacci, Tribonacci, and Ninbonacci---derived from generalized Fibonacci sequences and their Pisano periods under modulo 9 arithmetic (24, 24, 24, 8, and 1 respectively), creating a 72-8-1 orbital structure. We then introduce a primitive/non-primitive dichotomy based on greatest common divisor analysis, followed by a classification of primitive triples into three classical subfamilies: Fermat ($|C-F|=1$), Pythagoras ($(d-e)^2=1$), and Plato ($|G-F|=2$). We further record a ``gender'' classification distinguishing male tuples (primitive or $\gcd \neq 2$) from female tuples ($\gcd=2$, octave harmonics). Finally, we identify conjectural directions relating this exact integer taxonomy to E8-style lattice features and quantum-arithmetic signal metrics. These latter connections are stated as validation targets, not as consequences of the classification theorem.
\end{abstract}

\section{Introduction}

\subsection{Background on Pythagorean Triples}

A \emph{Pythagorean triple} is an ordered triple of positive integers $(C, F, G)$ satisfying the Pythagorean theorem:
\begin{equation}
C^2 + F^2 = G^2
\end{equation}

The classical parametric representation of primitive Pythagorean triples, known since Euclid, is given by:
\begin{equation}
C = m^2 - n^2, \quad F = 2mn, \quad G = m^2 + n^2
\end{equation}
where $m > n > 0$, $\gcd(m,n) = 1$, and $m \not\equiv n \pmod{2}$.

While this parametrization is well-established, it does not reveal deeper structural patterns in the distribution of Pythagorean triples. In this paper, we introduce an alternative generating framework based on generalized Fibonacci sequences that reveals a hidden five-fold symmetry.

\subsection{The BEDA Framework}

We introduce the \emph{BEDA tuple} representation, defined by four parameters $(b, e, d, a)$ satisfying:
\begin{align}
d &= b + e \label{eq:beda1}\\
a &= b + 2e \label{eq:beda2}
\end{align}

From any BEDA tuple, we generate a Pythagorean triple via:
\begin{align}
C &= 2de \label{eq:pyth1}\\
F &= ab \label{eq:pyth2}\\
G &= e^2 + d^2 \label{eq:pyth3}
\end{align}

\begin{proposition}
Every BEDA tuple $(b, e, d, a)$ satisfying equations \eqref{eq:beda1}--\eqref{eq:beda2} generates a valid Pythagorean triple via equations \eqref{eq:pyth1}--\eqref{eq:pyth3}.
\end{proposition}

\begin{proof}
We verify $C^2 + F^2 = G^2$:
\begin{align*}
C^2 + F^2 &= (2de)^2 + (ab)^2 \\
&= 4d^2e^2 + a^2b^2 \\
&= 4(b+e)^2e^2 + (b+2e)^2b^2 \\
&= 4(b^2e^2 + 2be^3 + e^4) + (b^4 + 4b^3e + 4b^2e^2) \\
&= 4b^2e^2 + 8be^3 + 4e^4 + b^4 + 4b^3e + 4b^2e^2 \\
&= b^4 + 4b^3e + 8b^2e^2 + 8be^3 + 4e^4
\end{align*}

Meanwhile:
\begin{align*}
G^2 &= (e^2 + d^2)^2 = (e^2 + (b+e)^2)^2 \\
&= (e^2 + b^2 + 2be + e^2)^2 = (b^2 + 2be + 2e^2)^2 \\
&= b^4 + 4b^3e + 8b^2e^2 + 8be^3 + 4e^4
\end{align*}

Therefore $C^2 + F^2 = G^2$.
\end{proof}

\subsection{Digital Roots and Pisano Periods}

The \emph{digital root} of a positive integer $n$, denoted $\text{dr}(n)$, is defined as:
\begin{equation}
\text{dr}(n) = \begin{cases}
9 & \text{if } n \equiv 0 \pmod{9} \\
n \bmod 9 & \text{otherwise}
\end{cases}
\end{equation}

The digital root maps integers to the set $\{1, 2, 3, 4, 5, 6, 7, 8, 9\}$.

For a Fibonacci-like sequence $\{F_n\}$ defined by the recurrence $F_n = F_{n-1} + F_{n-2}$ with initial conditions $F_0 = a, F_1 = b$, the \emph{Pisano period} $\pi(m)$ is the period with which the sequence repeats when taken modulo $m$. In this work, we focus specifically on $m = 9$.

\section{The Five Families}

\begin{definition}[The Five Families]
We define five generalized Fibonacci sequences, each with initial terms $(F_0, F_1)$ and recurrence relation $F_n = F_{n-1} + F_{n-2}$:

\begin{enumerate}
\item \textbf{Fibonacci Family}: $(F_0, F_1) = (1, 1)$
\item \textbf{Lucas Family}: $(F_0, F_1) = (2, 1)$
\item \textbf{Phibonacci Family}: $(F_0, F_1) = (3, 1)$
\item \textbf{Tribonacci Family}: $(F_0, F_1) = (3, 3)$
\item \textbf{Ninbonacci Family}: $(F_0, F_1) = (9, 9)$
\end{enumerate}
\end{definition}

Each family generates a sequence of $(b, e)$ pairs via $(b, e) = (F_n, F_{n+1})$.

\subsection{Pisano Periods of the Five Families}

\begin{theorem}[Pisano Periods under Modulo 9]\label{thm:pisano}
The five families exhibit the following Pisano periods when their digital roots are computed modulo 9:

\begin{center}
\begin{tabular}{lcc}
\toprule
\textbf{Family} & \textbf{Initial Terms} & \textbf{Pisano Period $\pi(9)$} \\
\midrule
Fibonacci & $(1, 1)$ & 24 \\
Lucas & $(2, 1)$ & 24 \\
Phibonacci & $(3, 1)$ & 24 \\
Tribonacci & $(3, 3)$ & 8 \\
Ninbonacci & $(9, 9)$ & 1 \\
\bottomrule
\end{tabular}
\end{center}
\end{theorem}

\begin{proof}
The Pisano periods are established through direct computation. For each sequence, we compute the digital roots of consecutive terms and detect when the pattern first repeats. The computations are provided in Appendix A, with full verification available in the accompanying code repository.

For example, the Fibonacci sequence modulo 9 produces the digital root cycle:
\[
\{1, 1, 2, 3, 5, 8, 4, 3, 7, 1, 8, 9, 8, 8, 7, 6, 4, 1, 5, 6, 2, 8, 1, 9\}
\]
which repeats after 24 terms.
\end{proof}

\section{Main Result: Complete Partition}

\begin{theorem}[Complete Partition of Digital Root Pairs]\label{thm:main}
The five families partition the set of all possible digital root pairs $\{1, 2, \ldots, 9\}^2$ into five disjoint subsets. Specifically:

\begin{enumerate}
\item The Fibonacci family contributes 24 unique $(\text{dr}(b), \text{dr}(e))$ pairs
\item The Lucas family contributes 24 unique $(\text{dr}(b), \text{dr}(e))$ pairs
\item The Phibonacci family contributes 24 unique $(\text{dr}(b), \text{dr}(e))$ pairs
\item The Tribonacci family contributes 8 unique $(\text{dr}(b), \text{dr}(e))$ pairs
\item The Ninbonacci family contributes 1 unique $(\text{dr}(b), \text{dr}(e))$ pair
\end{enumerate}

The total is $24 + 24 + 24 + 8 + 1 = 81 = 9^2$, accounting for all possible digital root pairs with \emph{no overlaps}.
\end{theorem}

\begin{proof}
We establish this result computationally by:
\begin{enumerate}
\item Computing the Pisano cycles for each of the five families
\item Extracting the unique $(\text{dr}(F_n), \text{dr}(F_{n+1}))$ pairs from each cycle
\item Verifying that the union of these five sets contains exactly 81 pairs
\item Verifying that the intersection of any two distinct sets is empty
\end{enumerate}

The complete classification table is provided in Table \ref{tab:classification}.
\end{proof}

\begin{table}[h!]
\centering
\caption{Complete Classification of Digital Root Pairs by Family}
\label{tab:classification}
\small
\begin{tabular}{c|ccccccccc}
\toprule
$\text{dr}(b) \backslash \text{dr}(e)$ & 1 & 2 & 3 & 4 & 5 & 6 & 7 & 8 & 9 \\
\midrule
1 & F & F & L & P & F & P & L & F & F \\
2 & L & L & F & L & P & P & L & F & L \\
3 & P & P & T & L & F & T & F & L & T \\
4 & F & P & F & P & P & L & L & P & P \\
5 & P & L & L & P & P & F & P & F & P \\
6 & L & F & T & F & L & T & P & P & T \\
7 & F & L & P & P & L & F & L & L & L \\
8 & F & L & P & F & P & L & F & F & F \\
9 & F & L & T & P & P & T & L & F & N \\
\bottomrule
\end{tabular}

\vspace{0.2cm}
\small{F = Fibonacci, L = Lucas, P = Phibonacci, T = Tribonacci, N = Ninbonacci}
\end{table}

\subsection{Primitive and Non-Primitive Triples}

While the five families provide a complete partition of digital root pairs, Pythagorean triples themselves admit a further fundamental dichotomy based on their greatest common divisor.

\begin{definition}[Primitive Pythagorean Triple]
A Pythagorean triple $(C, F, G)$ is \emph{primitive} if $\gcd(C, F, G) = 1$. Otherwise, it is \emph{non-primitive}.
\end{definition}

Every non-primitive triple is a positive integer multiple of a unique primitive triple. For example, $(6, 8, 10) = 2 \cdot (3, 4, 5)$ and $(9, 12, 15) = 3 \cdot (3, 4, 5)$ are both non-primitive, derived from the primitive triple $(3, 4, 5)$.

\begin{proposition}[Scaling and Family Labels]
Primitive parentage and digital-root family membership are distinct classifications. Scaling a primitive BEDA tuple by $k$ preserves the underlying similarity class of the Pythagorean triple, producing a $k^2$ multiple of the triple, but it need not preserve the named digital-root family label.
\end{proposition}

\begin{proof}
The scaled BEDA tuple $(kb,ke,kd,ka)$ generates $k^2(C,F,G)$, so it has the same primitive parent after dividing by the common factor. However, the digital-root pair changes from $(\mathrm{dr}(b),\mathrm{dr}(e))$ to $(\mathrm{dr}(kb),\mathrm{dr}(ke))$, which may lie in a different named family. For example, $(b,e)=(1,1)$ is Fibonacci and gives $(C,F,G)=(4,3,5)$, while its double $(2,2)$ gives $(16,12,20)=4(4,3,5)$ but has digital-root pair $(2,2)$, a Lucas cell in Table~\ref{tab:classification}. Thus the primitive parent is preserved by BEDA scaling, but the named digital-root family label is not.
\end{proof}

This distinction is important: the five-family partition classifies BEDA residue pairs, while primitivity classifies the Pythagorean triple itself. Non-primitives should therefore record both their current digital-root family and their primitive parent.

\begin{table}[h!]
\centering
\caption{Examples of Primitive and Non-Primitive Triples by Family}
\label{tab:primitive_examples}
\small
\begin{tabular}{lcccc}
\toprule
\textbf{Family} & \textbf{Primitive example} & \textbf{Primitive $(C,F,G)$} & \textbf{Non-primitive example} & \textbf{Parent} \\
\midrule
Fibonacci & $(1,1)$ & $(4,3,5)$ & $(2,3)\mapsto(30,16,34)$ & $2(15,8,17)$ \\
Lucas & $(1,3)$ & $(24,7,25)$ & $(2,2)\mapsto(16,12,20)$ & $4(4,3,5)$ \\
Phibonacci & $(1,4)$ & $(40,9,41)$ & $(4,2)\mapsto(24,32,40)$ & $8(3,4,5)$ \\
Tribonacci & none in primitive layer & --- & $(3,3)\mapsto(36,27,45)$ & $9(4,3,5)$ \\
Ninbonacci & none in primitive layer & --- & $(9,9)\mapsto(324,243,405)$ & $81(4,3,5)$ \\
\bottomrule
\end{tabular}
\end{table}

\section{Implications and Orbital Structure}

\subsection{The 72-8-1 Structure}

The five families naturally group into three \emph{orbital classes} based on their Pisano periods:

\begin{itemize}
\item \textbf{24-Cycle Orbit ("Cosmos")}: Fibonacci + Lucas + Phibonacci = $24 + 24 + 24 = 72$ pairs
\item \textbf{8-Cycle Orbit ("Satellite")}: Tribonacci = 8 pairs
\item \textbf{1-Cycle Orbit ("Singularity")}: Ninbonacci = 1 pair (the fixed point $(9,9)$)
\end{itemize}

This 72-8-1 structure has been independently verified through computational experiments with Markovian dynamical systems on the BEDA parameter space.

\subsection{Connection to Pythagorean Triples}

\begin{corollary}
Every Pythagorean triple $(C, F, G)$ generated from a BEDA tuple $(b, e, d, a)$ can be uniquely classified into one of the five families based on $(\text{dr}(b), \text{dr}(e))$.
\end{corollary}

This provides a new perspective on the structure of Pythagorean triples, complementary to the classical $(m, n)$ parametrization.

\section{Classical Subfamilies of Primitive Triples}

Within the space of primitive Pythagorean triples, classical number theory has long recognized certain special families characterized by simple relationships among the triple components. We now show how these classical subfamilies admit elegant characterizations in the BEDA framework.

\subsection{The Three Classical Subfamilies}

\begin{definition}[Classical Subfamilies]\label{def:classical}
We define three subfamilies of primitive Pythagorean triples:

\begin{enumerate}
\item \textbf{Fermat Family}: Triples $(C, F, G)$ satisfying $|C - F| = 1$ (consecutive legs)
\item \textbf{Pythagoras Family}: Triples $(C, F, G)$ where the BEDA parameters satisfy $(d - e)^2 = 1$
\item \textbf{Plato Family}: Triples $(C, F, G)$ satisfying $|G - F| = 2$ (hypotenuse 2 more than one leg)
\end{enumerate}
\end{definition}

These three families capture different geometric properties of primitive Pythagorean triples and have been studied since antiquity.

\subsection{BEDA Characterizations}

The BEDA framework reveals elegant algebraic characterizations of these classical subfamilies:

\begin{theorem}[BEDA Characterizations of Classical Subfamilies]\label{thm:classical_beda}
Let $(C, F, G)$ be a primitive Pythagorean triple generated from BEDA tuple $(b, e, d, a)$ via equations \eqref{eq:pyth1}--\eqref{eq:pyth3}. Then:

\begin{enumerate}
\item $(C, F, G) \in$ \textbf{Fermat} $\iff |C - F| = |2e^2-b^2| = |b^2-2e^2| = 1$
\item $(C, F, G) \in$ \textbf{Pythagoras} $\iff (d - e)^2 = 1 \iff b = 1$
\item $(C, F, G) \in$ \textbf{Plato} $\iff |G - F| = 2 \iff e = 1$; with the primitive restriction this is $e=1$ and $b$ odd
\end{enumerate}
\end{theorem}

\begin{proof}
These characterizations follow directly from the BEDA-to-Pythagorean transformation \eqref{eq:pyth1}--\eqref{eq:pyth3}. For Fermat,
\[
C-F = 2de-ab = 2e(b+e)-b(b+2e)=2e^2-b^2,
\]
so $|C-F|=1$ is the Pell-boundary condition $|b^2-2e^2|=1$. The Pythagoras family condition $(d-e)^2 = 1$ is particularly elegant: since $d = b + e$, this becomes $(b+e-e)^2 = b^2 = 1$, which occurs when $b = 1$ for positive BEDA tuples. Finally,
\[
G-F=(d^2+e^2)-(d^2-e^2)=2e^2,
\]
so $|G-F|=2$ occurs exactly when $e=1$; imposing primitivity further requires $d-e=b$ odd.
\end{proof}

\subsection{Examples of Classical Subfamilies}

\begin{table}[h!]
\centering
\caption{Examples of Classical Subfamilies Across the Five Families}
\label{tab:classical_examples}
\small
\begin{tabular}{llccc}
\toprule
\textbf{Subfamily} & \textbf{Five Family} & \textbf{$(b,e)$} & \textbf{$(C,F,G)$} & \textbf{Property} \\
\midrule
Fermat & Fibonacci & $(1,1)$ & $(4,3,5)$ & $4-3=1$ \\
Fermat & Phibonacci & $(3,2)$ & $(20,21,29)$ & $|20-21|=1$ \\
Fermat & Lucas & $(7,5)$ & $(120,119,169)$ & $120-119=1$ \\
Pythagoras & Fibonacci & $(1,1)$ & $(4,3,5)$ & $(d-e)^2=(2-1)^2=1$ \\
Pythagoras & Fibonacci & $(1,2)$ & $(12,5,13)$ & $(d-e)^2=(3-2)^2=1$ \\
Pythagoras & Lucas & $(1,3)$ & $(24,7,25)$ & $(d-e)^2=(4-3)^2=1$ \\
Plato & Fibonacci & $(1,1)$ & $(4,3,5)$ & $5-3=2$ \\
Plato & Phibonacci & $(3,1)$ & $(8,15,17)$ & $17-15=2$ \\
Plato & Fibonacci & $(7,1)$ & $(16,63,65)$ & $65-63=2$ \\
\bottomrule
\end{tabular}
\end{table}

Note that the classical subfamilies \emph{cut across} the five families, but not uniformly across all five in small bounded domains. In the examples above, Fermat appears in Fibonacci, Phibonacci, and Lucas classes; Pythagoras appears in Fibonacci and Lucas classes; primitive Plato appears in Fibonacci and Phibonacci classes. This demonstrates that the five-family partition (Layer 0) and the classical subfamily structure (Layer 2) are distinct classification schemes.

\begin{remark}[The Elegance of $(d-e)^2=1$]
The Pythagoras family characterization $(d-e)^2 = 1$ provides a beautiful geometric insight: these triples correspond to BEDA tuples lying exactly one unit off the diagonal $d = e$ in the integer lattice. This ``1-step-off-diagonal'' property makes them particularly amenable to computational generation and analysis.
\end{remark}

\section{Examples}

We provide examples of Pythagorean triples from each family:

\subsection{Fibonacci Family}
From $(b, e) = (1, 1)$: $(C, F, G) = (4, 3, 5)$\\
From $(b, e) = (2, 3)$: $(C, F, G) = (30, 16, 34)$\\
From $(b, e) = (5, 8)$: $(C, F, G) = (208, 105, 233)$

\subsection{Lucas Family}
From $(b, e) = (2, 1)$: $(C, F, G) = (6, 8, 10)$\\
From $(b, e) = (3, 4)$: $(C, F, G) = (56, 33, 65)$\\
From $(b, e) = (7, 11)$: $(C, F, G) = (396, 203, 445)$

\subsection{Tribonacci Family (8-Cycle)}
From $(b, e) = (3, 3)$: $(C, F, G) = (36, 27, 45)$\\
From $(b, e) = (6, 9)$: $(C, F, G) = (270, 144, 306)$\\
From $(b, e) = (15, 24)$: $(C, F, G) = (1872, 945, 2097)$

\subsection{Phibonacci Family}
From $(b, e) = (3, 1)$: $(C, F, G) = (8, 15, 17)$\\
From $(b, e) = (4, 5)$: $(C, F, G) = (90, 56, 106)$\\
From $(b, e) = (9, 14)$: $(C, F, G) = (644, 333, 725)$

\subsection{Ninbonacci Family (Fixed Point)}
From $(b, e) = (9, 9)$: $(C, F, G) = (324, 243, 405)$\\
All multiples: $(648, 486, 810)$, $(972, 729, 1215)$, etc.

\section{Gender Classification and Octave Harmonics}

A remarkable pattern emerges when we examine non-primitive triples with $\gcd = 2$: they exhibit a systematic ``octave'' relationship with primitive triples, analogous to harmonic doubling in music and signal processing.

\subsection{Male and Female Tuples}

\begin{definition}[Gender Classification]
For a Pythagorean triple $(C, F, G)$ with $\gcd(C, F, G) = k$, we define:

\begin{enumerate}
\item \textbf{Male tuple}: Primitive ($k = 1$) or non-primitive with $k \neq 2$
\item \textbf{Female tuple}: Non-primitive with $k = 2$ (exactly)
\item \textbf{Composite male}: Non-primitive with $k > 2$ and $k$ has odd prime factors
\end{enumerate}
\end{definition}

The terminology reflects a generative relationship: each primitive male tuple generates a unique female tuple via octave doubling.

\subsection{The Octave Transformation}

\begin{theorem}[Octave Transformation]\label{thm:octave}
Let $(b, e, d, a)$ be a BEDA tuple generating a primitive Pythagorean triple $(C, F, G)$ (a male tuple). Then the transformed BEDA tuple:
\begin{equation}
(b', e', d', a') = (2e, b, a, 2d) = (2e, b, b+2e, 2b+2e)
\end{equation}
generates a female tuple $(C', F', G')$ with $\gcd(C', F', G') = 2$.
\end{theorem}

\begin{proof}
Under the transformation $(b, e, d, a) \to (2e, b, a, 2d)$, we have:
\begin{align*}
d' &= b' + e' = 2e + b = a\\
a' &= b' + 2e' = 2e + 2b = 2d
\end{align*}
confirming the BEDA structure. The generated triple is:
\begin{align*}
C' &= 2d'e' = 2ab\\
F' &= a'b' = (2d)(2e) = 4de = 2(2de) = 2C\\
G' &= (e')^2 + (d')^2 = b^2 + a^2
\end{align*}

Since the original triple $(C, F, G) = (2de, ab, e^2 + d^2)$ was primitive, $\gcd(C, F, G) = 1$. The female triple is
\[
(C',F',G')=(2ab,4de,b^2+a^2)=(2F,2C,2G),
\]
because $b^2+a^2=b^2+(b+2e)^2=2(d^2+e^2)=2G$. Hence $\gcd(C',F',G')=2$ exactly. The transformation produces the first octave multiple of the primitive triple, with the two legs interchanged.
\end{proof}

\begin{example}[Male-Female Pairs]
Consider the primitive male tuple $(b,e) = (1,1)$ generating $(C,F,G) = (4,3,5)$. The female transformation gives:
\begin{align*}
(b', e', d', a') &= (2 \cdot 1, 1, 3, 2 \cdot 2) = (2, 1, 3, 4)\\
(C', F', G') &= (2 \cdot 3 \cdot 1, 4 \cdot 2, 1^2 + 3^2) = (6, 8, 10)
\end{align*}
Indeed, $(6,8,10) = 2 \cdot (3,4,5)$, confirming the octave relationship.
\end{example}

\subsection{Harmonic Interpretation}

The gender classification admits a simple harmonic interpretation:

\begin{itemize}
\item \textbf{Male tuples} represent fundamental frequencies
\item \textbf{Female tuples} represent first octave harmonics (frequency doubling)
\item \textbf{Composite male} tuples represent higher-order harmonics with complex factorization
\end{itemize}

This harmonic structure emerges naturally from the BEDA framework. Its usefulness as a signal-processing feature remains an empirical question.

\begin{table}[h!]
\centering
\caption{Male Primitives and Their Female Octaves}
\label{tab:gender}
\small
\begin{tabular}{lcccc}
\toprule
\textbf{Family} & \textbf{Male $(b,e)$} & \textbf{Male $(C,F,G)$} & \textbf{Female $(b',e')$} & \textbf{Female $(C',F',G')$} \\
\midrule
Fibonacci & $(1,1)$ & $(4,3,5)$ & $(2,1)$ & $(6,8,10)$ \\
Phibonacci & $(3,1)$ & $(8,15,17)$ & $(2,3)$ & $(30,16,34)$ \\
Lucas & $(1,2)$ & $(12,5,13)$ & $(4,1)$ & $(10,24,26)$ \\
Lucas & $(1,3)$ & $(24,7,25)$ & $(6,1)$ & $(14,48,50)$ \\
\bottomrule
\end{tabular}
\end{table}

\section{Conjectural E8 and Harmonicity Directions}

The exact classifications above suggest possible comparisons with highly symmetric lattice structures. This section records those comparisons as conjectural directions. The present paper does not prove an embedding of the BEDA taxonomy into the E8 root system.

\subsection{Background: The E8 Root System}

The E8 Lie algebra is an 8-dimensional exceptional simple Lie algebra with a root system consisting of 240 vectors in $\mathbb{R}^8$ exhibiting extraordinary symmetry. The E8 root system can be constructed via several polytopes, including the Gosset $4_{21}$ polytope, and has been extensively studied in mathematics and theoretical physics.

The E8 lattice admits a weight lattice structure with concentric shells:
\begin{itemize}
\item \textbf{Root shell}: 240 shortest vectors (the E8 roots themselves)
\item \textbf{First weight shell}: 2160 vectors at distance $\sqrt{2} \times$ root length
\item \textbf{Higher weight shells}: Additional concentric shells at increasing distances
\end{itemize}

\subsection{BEDA Feature Embedding}

While BEDA tuples are 4-dimensional $(b, e, d, a)$, they can be doubled into an 8-dimensional feature vector via the mapping:
\begin{equation}
\phi: (b, e, d, a) \mapsto (b, e, d, a, b, e, d, a) \in \mathbb{R}^8
\end{equation}

This map is a feature embedding only. By itself it is not an isometry into the E8 lattice, does not prove E8 shell membership, and does not identify BEDA tuples with E8 roots.

\begin{conjecture}[E8 Alignment Program]\label{conj:e8}
Let $(C, F, G)$ be a Pythagorean triple with $\gcd(C, F, G) = k$ generated from BEDA tuple $(b, e, d, a)$. There may exist a normalization and a lattice embedding, extending or replacing the feature map $\phi$, under which the following comparisons become meaningful test targets:

\begin{enumerate}
\item \textbf{Primitive male tuples} ($k=1$): candidate alignment with a primitive shell or root-like layer
\item \textbf{Female tuples} ($k=2$): candidate alignment with a first doubled or octave shell
\item \textbf{Composite male tuples} ($k>2$): candidate alignment with higher scaled shells
\end{enumerate}
\end{conjecture}

\paragraph{Validation requirements.}
Conjecture~\ref{conj:e8} is not proved here. A proof or disproof would require an explicit normalization, a specified lattice map, exact shell-membership predicates, and counterexample searches over BEDA-generated triples. The simple scaling identity $(b,e,d,a) \mapsto (kb,ke,kd,ka)$ changes Euclidean feature-vector length by $k$ and changes the associated triple by $k^2$, so scaling alone is not evidence for E8 shell membership.

\subsection{Interpretive Hypotheses for Geometric Harmonicity}

If a valid lattice embedding is later found, it could provide a geometric interpretation of the gender classification:

\begin{itemize}
\item \textbf{Primitive males} would serve as candidate fundamental geometric units
\item \textbf{Females} would mark a systematic octave relation to primitive parents
\item \textbf{Composite males} would continue the scaling hierarchy into higher layers
\end{itemize}

These are hypotheses for future work. The proven content in this paper is the integer taxonomy and the BEDA formulas, not an E8 inheritance theorem.

\begin{remark}[Connection to Quantum Arithmetic]
In quantum arithmetic (QA) systems, the BEDA taxonomy can be used as an exact source of modular and harmonic features. Whether these features improve signal classifiers, lattice-quality metrics, or other empirical systems is a separate experimental question.
\end{remark}

\begin{definition}[Candidate Harmonicity Index]
For a Pythagorean triple $(C, F, G)$ generated from BEDA tuple $(b, e, d, a)$, one possible \emph{Harmonicity Index} feature score is:
\begin{equation}
HI_2 = w_{\text{ang}} \cdot H_{\text{angular}} + w_{\text{rad}} \cdot H_{\text{radial}} + w_{\text{fam}} \cdot H_{\text{family}}
\end{equation}
where:
\begin{itemize}
\item $H_{\text{angular}}$: Angular harmonicity based on $(\text{mod}\, 24) \times (\text{mod}\, 9)$ structure
\item $H_{\text{radial}} = 1/\gcd(C,F,G)$: Radial harmonicity (primitivity measure)
\item $H_{\text{family}}$: Membership in Fermat/Pythagoras/Plato classical subfamilies
\end{itemize}

This is a proposed feature bundle. No claim is made here that it encodes E8 shell membership or predicts signal behavior without separate empirical validation.
\end{definition}

\section{Potential Applications and Validation Targets}

\subsection{Computational Number Theory}

The classification provides exact algorithms for:
\begin{itemize}
\item Generating Pythagorean triples by family
\item Classifying existing triples into families
\item Filtering and summarizing large Pythagorean-triple datasets via digital-root analysis
\end{itemize}

\subsection{Dynamical Systems}

The 72-8-1 orbital structure arises naturally in discrete dynamical systems governed by modular arithmetic. This connects the static classification of Pythagorean triples to the dynamics of recurrence relations.

\subsection{Quantum Arithmetic and Signal Processing}

The hierarchical classification presented here supplies exact feature labels that can be tested inside QA signal-processing experiments. Candidate features include:

\begin{itemize}
\item \textbf{Harmonicity Index 2.0}: a feature score based on primitivity, gender, and classical subfamily structure
\item \textbf{Lattice-style alignment scores}: optional geometric features requiring independent validation
\item \textbf{Family-aware divergence measures}: candidate comparison metrics for empirical datasets
\item \textbf{Fermat/Pythagoras/Plato subfamily membership}: exact labels whose empirical relevance must be measured case by case
\end{itemize}

No empirical performance claim is made in this paper.

\subsection{Cryptographic Caution}

The modular orbit structure may be useful for constructing toy protocols or test problems. However, this paper makes no cryptographic security claim, proposes no trapdoor function, and does not reduce any family-classification problem to a standard hard problem.

\section{Conclusion and Future Work}

We have presented a comprehensive hierarchical taxonomy of Pythagorean triples organized into three nested layers:

\begin{enumerate}
\item \textbf{Layer 0: Five Families} (Fibonacci, Lucas, Phibonacci, Tribonacci, Ninbonacci) partitioning all 81 digital root pairs via Pisano periods 24-24-24-8-1, creating a 72-8-1 orbital structure
\item \textbf{Layer 1: Primitive/Non-Primitive Dichotomy} based on $\gcd(C,F,G)$, distinguishing primitive parentage from the current digital-root family label of a BEDA tuple
\item \textbf{Layer 2: Classical Subfamilies} (Fermat, Pythagoras, Plato) cutting across the five families with elegant BEDA characterizations
\end{enumerate}

Beyond this static classification, we introduced a \textbf{gender classification} distinguishing male tuples (primitive or $\gcd \neq 2$), female tuples ($\gcd = 2$, octave harmonics), and composite males ($\gcd > 2$), revealing systematic generative relationships via the octave transformation $(b,e,d,a) \to (2e, b, a, 2d)$.

We also identified a conjectural \textbf{E8 alignment program}: the BEDA feature embedding and octave hierarchy suggest possible comparisons with lattice shells, but a genuine E8 correspondence requires a separate exact embedding theorem and counterexample search.

The unified framework connects classical number theory, modular arithmetic, Fibonacci sequences, and harmonic integer features. Proposed links to lattice geometry, QA signal processing, and cryptographic constructions remain future validation targets.

Future research directions include:
\begin{itemize}
\item \textbf{Higher moduli}: Extending Pisano period analysis to $m > 9$ and discovering additional family structures
\item \textbf{E8 embedding tests}: Developing and falsifying explicit candidate maps between BEDA features and E8-style lattice shells
\item \textbf{Modular forms connection}: Investigating relationships between the 72-8-1 structure and modular form weight spaces
\item \textbf{Security analysis}: Testing whether any family-based toy protocol has a credible hardness assumption
\item \textbf{Quantum arithmetic applications}: Testing HI 2.0 and related exact features in controlled signal-classification experiments
\item \textbf{Higher-dimensional generalizations}: Extending BEDA framework to Pythagorean quadruples, quintuples, and beyond
\item \textbf{Computational verification}: Exhaustive exact testing of candidate lattice embeddings and family refinements over large BEDA domains
\end{itemize}

\section*{Acknowledgments}

[To be added]

\begin{thebibliography}{99}

\bibitem{euclid}
Euclid, \emph{Elements}, Book X.

\bibitem{dickson}
L. E. Dickson, \emph{History of the Theory of Numbers, Vol. II: Diophantine Analysis}, Carnegie Institution of Washington, 1920.

\bibitem{hardy}
G. H. Hardy and E. M. Wright, \emph{An Introduction to the Theory of Numbers}, Oxford University Press, 6th edition, 2008.

\bibitem{pisano}
D. D. Wall, ``Fibonacci Series Modulo $m$'', \emph{American Mathematical Monthly}, Vol. 67, No. 6 (1960), pp. 525-532.

\bibitem{humphreys}
J. E. Humphreys, \emph{Introduction to Lie Algebras and Representation Theory}, Springer-Verlag, 1972.

\bibitem{conway_sloane}
J. H. Conway and N. J. A. Sloane, \emph{Sphere Packings, Lattices and Groups}, Springer-Verlag, 3rd edition, 1999.

\bibitem{baez_e8}
J. C. Baez, ``The Octonions'', \emph{Bulletin of the American Mathematical Society}, Vol. 39, No. 2 (2002), pp. 145-205.

\bibitem{coxeter}
H. S. M. Coxeter, ``Integral Cayley Numbers'', \emph{Duke Mathematical Journal}, Vol. 13, No. 4 (1946), pp. 561-578.

\bibitem{qa_foundations}
[QA Research Team], ``Quantum Arithmetic Foundations and E8 Alignment Metrics'', \emph{In preparation}, 2025.

\end{thebibliography}

\appendix

\section{Computational Verification}

The arithmetic examples and finite-domain checks used in preparing this draft were verified using Python 3.x exact integer arithmetic. The complete source code, including:
\begin{itemize}
\item Generators for all five families
\item Pisano period computation
\item Partition verification
\item Pythagorean triple generation and validation
\end{itemize}

is available at: \url{https://github.com/1r0nw1ll/quantum-arithmetic-research}

\noindent The four main theorems of this paper have independent machine-checkable certificate validators in the repository's QA certificate system:
\begin{itemize}
\item \textbf{Theorem \ref{thm:pisano}} (Pisano Periods): \texttt{qa\_alphageometry\_ptolemy/qa\_pisano\_orbit\_correspondence\_cert\_v1/}
\item \textbf{Theorem \ref{thm:main}} (Complete Partition): \texttt{qa\_alphageometry\_ptolemy/qa\_five\_families\_partition\_cert\_v1/}
\item \textbf{Theorem \ref{thm:classical_beda}} (BEDA Classical Subfamilies): \texttt{qa\_alphageometry\_ptolemy/qa\_beda\_classical\_subfamilies\_cert\_v1/}
\item \textbf{Theorem \ref{thm:octave}} (Octave Transformation): \texttt{qa\_alphageometry\_ptolemy/qa\_octave\_transformation\_cert\_v1/}
\end{itemize}

Each validator runs self-contained Python 3 exact integer arithmetic, produces a JSON result with \texttt{"ok": true}, and has a pinned SHA-256 hash for reproducibility.

\section{Generated Classical Slice Tables}

The following appendix tables were generated from the exact BEDA laws using the script \texttt{tools/generate\_pythagorean\_classical\_tables.py}. The companion JSON evidence file is \texttt{results/pythagorean\_classical\_tables.json}.

% Generated by tools/generate_pythagorean_classical_tables.py
\begin{table}[h!]
\centering
\caption{Generated Fermat/Pell Boundary Rows}
\label{tab:appendix_fermat_pell}
\small
\begin{tabular}{llcccl}
\toprule
\textbf{$(b,e)$} & \textbf{Family} & \textbf{$(C,F,G)$} & \textbf{$\gcd$} & \textbf{Primitive} & \textbf{Check} \\
\midrule
$(1,1)$ & Fibonacci & $(4,3,5)$ & 1 & yes & $b^2-2e^2=-1$ \\
$(3,2)$ & Phibonacci & $(20,21,29)$ & 1 & yes & $b^2-2e^2=1$ \\
$(7,5)$ & Lucas & $(120,119,169)$ & 1 & yes & $b^2-2e^2=-1$ \\
$(17,12)$ & Phibonacci & $(696,697,985)$ & 1 & yes & $b^2-2e^2=1$ \\
$(41,29)$ & Lucas & $(4060,4059,5741)$ & 1 & yes & $b^2-2e^2=-1$ \\
$(99,70)$ & Lucas & $(23660,23661,33461)$ & 1 & yes & $b^2-2e^2=1$ \\
$(239,169)$ & Phibonacci & $(137904,137903,195025)$ & 1 & yes & $b^2-2e^2=-1$ \\
$(577,408)$ & Lucas & $(803760,803761,1136689)$ & 1 & yes & $b^2-2e^2=1$ \\
\bottomrule
\end{tabular}
\end{table}

\begin{table}[h!]
\centering
\caption{Generated Pythagoras Slice Rows ($b=1$)}
\label{tab:appendix_pythagoras}
\small
\begin{tabular}{llcccl}
\toprule
\textbf{$(b,e)$} & \textbf{Family} & \textbf{$(C,F,G)$} & \textbf{$\gcd$} & \textbf{Primitive} & \textbf{Check} \\
\midrule
$(1,1)$ & Fibonacci & $(4,3,5)$ & 1 & yes & $b=1$ \\
$(1,2)$ & Fibonacci & $(12,5,13)$ & 1 & yes & $b=1$ \\
$(1,3)$ & Lucas & $(24,7,25)$ & 1 & yes & $b=1$ \\
$(1,4)$ & Phibonacci & $(40,9,41)$ & 1 & yes & $b=1$ \\
$(1,5)$ & Fibonacci & $(60,11,61)$ & 1 & yes & $b=1$ \\
$(1,6)$ & Phibonacci & $(84,13,85)$ & 1 & yes & $b=1$ \\
$(1,7)$ & Lucas & $(112,15,113)$ & 1 & yes & $b=1$ \\
$(1,8)$ & Fibonacci & $(144,17,145)$ & 1 & yes & $b=1$ \\
\bottomrule
\end{tabular}
\end{table}

\begin{table}[h!]
\centering
\caption{Generated Plato Slice Rows ($e=1$)}
\label{tab:appendix_plato}
\small
\begin{tabular}{llcccl}
\toprule
\textbf{$(b,e)$} & \textbf{Family} & \textbf{$(C,F,G)$} & \textbf{$\gcd$} & \textbf{Primitive} & \textbf{Check} \\
\midrule
$(1,1)$ & Fibonacci & $(4,3,5)$ & 1 & yes & $e=1$ \\
$(2,1)$ & Lucas & $(6,8,10)$ & 2 & no & $e=1$ \\
$(3,1)$ & Phibonacci & $(8,15,17)$ & 1 & yes & $e=1$ \\
$(4,1)$ & Fibonacci & $(10,24,26)$ & 2 & no & $e=1$ \\
$(5,1)$ & Phibonacci & $(12,35,37)$ & 1 & yes & $e=1$ \\
$(6,1)$ & Lucas & $(14,48,50)$ & 2 & no & $e=1$ \\
$(7,1)$ & Fibonacci & $(16,63,65)$ & 1 & yes & $e=1$ \\
$(8,1)$ & Fibonacci & $(18,80,82)$ & 2 & no & $e=1$ \\
$(9,1)$ & Fibonacci & $(20,99,101)$ & 1 & yes & $e=1$ \\
$(10,1)$ & Fibonacci & $(22,120,122)$ & 2 & no & $e=1$ \\
\bottomrule
\end{tabular}
\end{table}

\end{document}
